\section{The Epistemic Super-Domain Layer and the Narrator}
\label{sec:epistemic}

This section describes the framework's distinctive analytical layer: the categorical epistemic typing on every relation, the deterministic rule engine that produces it, the per-domain customization paths, and the persona-driven \emph{narrator} stage that synthesizes the classified graph into an analyst briefing. Three output artifacts are produced together, and each has its own consumer:

\begin{itemize}
\item \texttt{graph\_data.json} -- Layer 1, the brute-facts graph (Section~\ref{sec:architecture}).
\item \texttt{claims\_layer.json} -- Layer 2, the deterministic per-relation epistemic classification produced by the rule engine. Reproducible bit-identically across runs given identical Layer 1 input. The canonical machine-readable record.
\item \texttt{epistemic\_narrative.md} -- the analyst briefing produced by the persona-driven narrator stage from the classified graph (\emph{not} from the source documents). Structured into Executive Summary, Asserted Claims, Prophetic Claims, Contested Claims, Gaps, and Recommended Follow-Ups. Non-deterministic in word choice but stable in structure and findings across runs.
\end{itemize}

A fourth surface, the \emph{workbench chat}, is the reactive consumer of the same classified graph and the same domain persona. Together, the classified graph (\texttt{claims\_layer.json}), the proactive briefing (\texttt{epistemic\_narrative.md}), and the reactive chat are what the framework adds beyond a typed-relation knowledge graph -- and they are what an analyst actually consumes when working with the framework on a corpus.

\subsection{Rationale}

A confidence score in $[0, 1]$ on a relation is a one-dimensional signal: it conveys how strongly the extractor believes a given subject--predicate--object triple is in the source text. It does not convey what \emph{kind of claim} the source text is making. The same edge in the graph could be supported by source language that asserts ($p<0.001$ in a randomized trial), prophesies (``in one or more embodiments\ldots'' in a patent), hypothesizes (``has been suggested to\ldots'' in a discussion section), or denies (``did not show a statistically significant\ldots''). These differ \emph{categorically}, not by degree, and a scalar confidence flattens them all into the same numerical surface.

The epistemic super-domain layer adds a second dimension. Every relation in the graph receives a categorical \emph{epistemic status} from a small fixed taxonomy. The taxonomy was chosen to be (a) source-language-driven (every status corresponds to identifiable patterns in the verbatim evidence text), (b) decision-relevant (analysts make different decisions about \texttt{prophetic} vs.\ \texttt{asserted} edges), and (c) cross-domain general (the same taxonomy applies in drug-discovery, contracts, clinicaltrials, and FDA labels, with per-domain variations).

\subsection{The status taxonomy}

The shared epistemic status taxonomy contains seven categorical labels (Table~\ref{tab:epistemic-statuses}). Six are positive classifications (a relation has been classified as one of these); the seventh, \texttt{unclassified}, is assigned when no rule matches. Domain-specific epistemic modules may add or refine statuses (for example, \texttt{fda-product-labels} introduces a four-level classifier mapping onto these statuses with finer granularity).

\begin{table}[t]
\centering
\small
\caption{The shared epistemic status taxonomy.}
\label{tab:epistemic-statuses}
\begin{tabular}{@{}lp{8.5cm}@{}}
\toprule
\textbf{Status} & \textbf{Meaning and trigger patterns} \\
\midrule
\texttt{asserted} & Quantitative or definitive claim grounded in the source. Default for evidence text containing trial results, p-values, dose-response data, or unhedged assertion. \\
\texttt{prophetic} & Patent forward-looking language. Triggered by phrases such as ``in one or more embodiments,'' ``may be useful for,'' ``the present invention provides,'' and document-type inference (the source document is a patent filing). \\
\texttt{hypothesized} & Hedged research language. Triggered by ``hypothesized,'' ``has been suggested,'' ``may indicate,'' ``is thought to,'' and similar epistemic hedging vocabulary. \\
\texttt{contested} & The same subject--predicate--object triple appears across multiple source documents with materially different confidence scores or contradictory evidence quotes. Detected at cross-source aggregation time, not within a single document. \\
\texttt{contradiction} & Explicit denial of a claim asserted in another source (``did not demonstrate,'' ``failed to show statistical significance,'' ``no evidence of\ldots''). Stronger than \texttt{contested}. \\
\texttt{negative} & Explicit absence claim from a single source (``no association observed,'' ``no detectable expression''). Distinct from \texttt{contradiction} (which is cross-source). \\
\texttt{speculative} & Marked future-conditional language outside patents (``further studies are needed,'' ``the implications remain to be determined''). \\
\texttt{unclassified} & No rule fired; default fall-through. \\
\bottomrule
\end{tabular}
\end{table}

\subsection{The rule engine}

Layer 2 is produced by a deterministic rule engine in \texttt{core/label\_epistemic.py}, parameterized by the domain's \texttt{epistemic.py} module. The engine has three stages, applied per relation:

\begin{enumerate}
\item \textbf{Document-type inference.} The source document's filename and any cached document-type hint (e.g., ``patent\_{}WO2024\_*\_zealand.txt'' triggers patent-type inference) drive a coarse classification. Patent-source relations default to \texttt{prophetic} unless a more specific trigger fires.
\item \textbf{Hedging-pattern matching.} A list of compiled regex patterns runs against the relation's verbatim evidence text. Each pattern carries a status it produces on match. Patterns are short, conservative, and audited: they prefer \texttt{unclassified} over a wrong classification.
\item \textbf{Cross-source aggregation.} After per-relation classification, the engine groups relations by $\langle$subject, predicate, object$\rangle$ identity and inspects the multi-source set. If the same triple appears with confidence ranges spanning more than 0.30 across sources, it is annotated \texttt{contested}. If a source's evidence text is a denial of another source's assertion, the deniers are annotated \texttt{contradiction}.
\end{enumerate}

The output of the rule engine is \texttt{claims\_layer.json}: a per-relation record with the assigned status, the matching rule (for audit), the verbatim evidence quote, and a back-reference to the source document. This artifact is bit-identical across runs given identical Layer 1 input. It is the reproducible Layer 2 record.

\subsection{Domain-specific epistemic rules}

Each domain's \texttt{epistemic.py} can add domain-specific rules beyond the shared taxonomy. Two examples illustrate the pattern.

\paragraph{Phase-based evidence grading (clinicaltrials).} Clinical trial documents from \texttt{ClinicalTrials.gov} carry trial-phase metadata (Phase 1, 2, 3, 4). The clinicaltrials epistemic module promotes Phase 3 and Phase 4 outcome relations to \texttt{asserted}, leaves Phase 2 at the rule-engine default, and demotes Phase 1 outcomes to \texttt{hypothesized} unless the evidence quote explicitly asserts a quantitative result. This phase-based grading captures regulatory-grade evidence stratification that a domain-agnostic rule engine cannot infer from text alone.

\paragraph{Four-level FDA epistemology (fda-product-labels).} FDA Structured Product Labels carry a recognizable epistemic vocabulary tied to label-section structure mandated by 21 CFR §201.57. The fda-product-labels epistemic module implements a four-level classifier mapped to that section structure with finer granularity than the shared taxonomy: \texttt{established} (definitive labeling language: BOXED WARNING, CONTRAINDICATION, mandated dosing limits), \texttt{observed} (controlled clinical study data from the CLINICAL STUDIES section: RCT-cited efficacy, dose-response, statistically significant outcomes), \texttt{reported} (post-market / spontaneous data from ADVERSE REACTIONS, POSTMARKETING, and DRUG INTERACTIONS sections: pharmacovigilance signals), and \texttt{theoretical} (mechanism-of-action inference from CLINICAL PHARMACOLOGY and PHARMACOLOGY sections: in-vitro and mechanistic reasoning). The four levels are visible to the analyst persona, which can stratify briefings accordingly.

These per-domain extensions illustrate the framework's pluggability: the core pipeline does not need to know about phase-based grading or FDA epistemology vocabulary -- the domain's \texttt{epistemic.py} encodes that knowledge, and \texttt{core/label\_epistemic.py} dispatches to it.

\subsection{Predefined defaults, customizable per domain}
\label{sec:epistemic-customization}

The framework ships with a usable epistemic layer out of the box. Every pre-built domain comes with: (a) the seven-status taxonomy in Table~\ref{tab:epistemic-statuses}; (b) a starter set of hedging-pattern regex rules that fire on common research-language constructions across all biomedical and contract corpora; (c) document-type inference defaults that recognize patent filings, trial-protocol records, FDA labels, and peer-reviewed papers from filename and structural cues; and (d) the cross-source aggregation rule that detects \texttt{contested} edges when the same triple appears with confidence variance above a threshold across documents. A user who installs the framework, picks a pre-built domain, and ingests their corpus gets all four behaviors without further configuration.

These defaults are the floor, not the ceiling. Three customization paths are supported:

\begin{itemize}
\item \emph{Add or refine hedging patterns.} The domain's \texttt{epistemic.py} exposes \texttt{HEDGING\_PATTERNS} as a list of $\langle$regex, status, rationale$\rangle$ triples. Adding a new pattern is one line; the rule engine picks it up on the next ingest.
\item \emph{Add domain-specific status types.} The clinicaltrials and fda-product-labels domains demonstrate how a domain can introduce statuses beyond the shared seven (\texttt{established}, \texttt{reported}, \texttt{theoretical} for FDA labels) by extending the rule engine. New domains can do the same; the rule-engine dispatch is open.
\item \emph{Override the cross-source aggregation threshold or add custom rules.} A domain can plug in \texttt{CUSTOM\_RULES} (cross-source contradiction patterns, domain-specific contested-pair detection, special-case escalations) without modifying \texttt{core/label\_epistemic.py}. The wizard generates a starter \texttt{epistemic.py} pre-populated with the shared taxonomy and the pattern of \texttt{HEDGING\_PATTERNS} / \texttt{CUSTOM\_RULES} so authors edit rather than write from scratch.
\end{itemize}

The customization story is part of the workspace pattern. A scientist exploring a new corpus type may discover a class of hedging language the defaults miss (e.g., contract-specific obligations stated in conditional-future tense, or laboratory-protocol-specific provisional language). The fix is a one-line addition to the domain's \texttt{epistemic.py}, not a feature request to the framework. This keeps the framework lean -- only the rules that matter to a given domain live in that domain -- and keeps domain authors in control of how their corpora are classified.

\subsection{The persona-driven narrator}

A separate stage consumes the classified graph and produces a proactive briefing in natural language. The narrator's system prompt is the domain's analyst persona, loaded directly from \texttt{workbench/template.yaml}. The input to the narrator is \emph{not} the source documents -- it is a structured summary of the classified graph, assembled from \texttt{claims\_layer.json}: per-status counts, contested-claim pairs with their confidence ranges, prophetic-claim concentrations by entity type, gap signals (entity types with low connection density), and any domain-specific signals the persona is configured to attend to. The narrator emits \texttt{epistemic\_narrative.md}, structured into Executive Summary, Asserted Claims, Prophetic Claims, Contested Claims, Gaps, and Recommended Follow-Ups sections.

A worked example: on the S6 GLP-1 corpus (drug-discovery domain, 34 documents), the narrator's executive summary identified \emph{61 prophetic claims} concentrated on cardiovascular, neurodegenerative, and metabolic-sub-disorder indications and recommended ``temporal stratification of confidence'' on the contested \texttt{semaglutide INDICATED\_FOR obesity} edge (confidence range $0.55$--$0.97$ across sources, where $0.55$ instances trace to pre-STEP-1 patent language and $0.97$ instances trace to post-approval clinical assertions). The narrator did not see the source documents directly; its claims trace back to graph nodes and relations annotated by the rule engine.

\subsection{Persona dual-use: one source, two surfaces}

The same persona string is also the system prompt for the workbench's reactive chat panel. When a user asks ``which prophetic claims are concentrated on cardiovascular outcomes?'' in chat, the answer is generated by the same persona that wrote the briefing, against the same classified graph, with the same vocabulary and decision criteria. This is by design: the proactive briefing and the reactive chat would otherwise drift apart in tone, vocabulary, and analytic framing if they were configured separately. The pattern reduces the surface area of the framework -- one YAML field defines the domain's analyst voice -- and makes domain extension cheaper: a domain author writes one persona, and both surfaces inherit it.
